\documentclass[12pt,a4paper]{article}
\usepackage[T1]{fontenc}
\usepackage[utf8]{inputenc}
\usepackage{amsmath,amssymb}
\usepackage{booktabs}
\usepackage{tabularx}
\usepackage{geometry}
\geometry{top=25mm, bottom=25mm, left=28mm, right=28mm}
\usepackage{setspace}
\onehalfspacing
\usepackage[colorlinks=true,linkcolor=blue!60!black,citecolor=blue!60!black]{hyperref}

\begin{document}

% ============================================================
% INSERT POINT: This section follows Section 3.17 (Summary) in
% Chapter 3: "Application to Hybrid Intrusion Detection on
% Network-Traffic Graphs and Experimental Validation"
% ============================================================

\section{Production Deployment and Platform Validation}
\label{sec:platform-validation}

The preceding sections of this chapter established the theoretical
foundations and empirical performance of seven mathematical methods
(M1--M7) for hybrid intrusion detection on network-traffic graphs.
Validation on benchmark datasets---including CICIDS-2017, NSL-KDD,
UNSW-NB15, and CIC-Bell-DNS-2021---demonstrated that spectral graph
wavelets (M1), persistent homology (M2), Finsler geometry (M3),
sheaf-theoretic consistency (M4), category-theoretic composition (M5),
information-geometric detection (M6), and tropical semiring
aggregation (M7) each contribute complementary detection capabilities
with quantifiable improvements over conventional baselines [143].
However, benchmark evaluation alone is insufficient to establish
operational readiness. This section extends validation to a
production-grade platform, RobustIDPS.ai (DOI:
10.5281/zenodo.19129512) [144], which integrates all seven methods
alongside additional innovations required for deployment within
operational Security Operations Centres (SOCs). The platform serves
both as a validation instrument for the dissertation's theoretical
contributions and as an independent research artefact demonstrating
the viability of mathematically grounded intrusion detection at
production scale.

% ────────────────────────────────────────────────────────────
\subsection{Platform Architecture and Scope}
\label{subsec:platform-architecture}

RobustIDPS.ai is a full-stack, cloud-native intrusion detection and
response platform designed to unify the research methods presented in
this dissertation with the operational requirements of contemporary
SOC workflows. The platform encompasses \textbf{12+ neural network
models} maintained in a unified model registry, exposed through
\textbf{46 interactive pages} organised across \textbf{9 functional
groups}, as summarised in Table~\ref{tab:functional-groups}.

The system architecture follows a layered microservice design.  The
backend is implemented in \textbf{FastAPI} (Python 3.11+), providing
asynchronous REST endpoints and WebSocket channels for real-time
streaming.  The frontend is a \textbf{React 18} single-page
application with responsive design and Progressive Web App (PWA)
manifest, enabling cross-platform accessibility on desktop, tablet,
and mobile devices without requiring native application installation
[145].  Persistent state is managed by \textbf{PostgreSQL 16} with
connection pooling, while \textbf{Redis 7} serves as both a caching
layer and a pub/sub broker for inter-service communication.  The
entire stack is containerised via \textbf{Docker Compose}, with
separate service definitions for the API server, the frontend build,
the database, the cache, and background worker processes.  Deployment
is reproducible through a single \texttt{docker compose up} invocation
[146].

A distinguishing architectural feature is the \textbf{Multi-LLM SOC
Copilot}, which supports simultaneous integration with Claude
(Anthropic), GPT-4o (OpenAI), Gemini (Google), and DeepSeek via a
provider-agnostic adapter layer.  This design permits SOC analysts to
select the most appropriate large language model for a given
investigative task---such as alert summarisation, natural-language
threat hunting, or incident report generation---while maintaining a
consistent conversational interface [147].

The \textbf{Live Monitor} page implements WebSocket-based streaming
with a configurable retention window of up to 2{,}000 flows, allowing
analysts to observe classification decisions in real time as network
traffic is ingested.  Each flow is processed through the selected
model pipeline and annotated with predicted class, confidence score,
and MITRE ATT\&CK technique mapping before being pushed to the
client.  A companion \textbf{File Replay} mode enables deterministic
re-ingestion of stored PCAP captures for reproducible validation
experiments.

\begin{table}[htbp]
\centering
\caption{RobustIDPS.ai functional groups and page inventory.}
\label{tab:functional-groups}
\small
\begin{tabularx}{\textwidth}{l c X}
\toprule
\textbf{Functional Group} & \textbf{Pages} & \textbf{Key Capabilities} \\
\midrule
AI Command Centre      & 5 & Dashboard, Upload \& Analyse, Analytics, Live Monitor, Alert Causality \\
AI Active Defence      & 5 & Red Team Arena, Threat Response, SOC Copilot, Federated Learning, Attack Chain Predictor \\
SOC Intelligence       & 6 & Auto-Investigation, Threat Hunt, Incident Reports, Threat Intel, Rule Generator, CVE Mapper \\
LLM Attack Surfaces    & 5 & Prompt Injection, Jailbreak Taxonomy, RAG Poisoning, Multi-Agent Chain, Data Poisoning \\
MLSecOps Standards     & 3 & MITRE ATT\&CK Mapper, Alert Triage, Compliance Hub (OWASP + NIST) \\
AI Security \& Gov     & 4 & PQ Cryptography, Zero-Trust, Supply Chain, Research Hub \\
AI Novel Methods       & 5 & Continual Learning, RL Response Agent, Adversarial Eval, Explainability, Autoencoder Detector \\
AI Data \& Models      & 3 & Datasets, Models (12+), Ablation Studio \\
Industry \& Research   & 4 & Lab Partnerships, Postdoc Portal, Interview Prep, Benchmarks \\
\midrule
\textbf{Total}         & \textbf{40+} & \\
\bottomrule
\end{tabularx}
\end{table}

% ────────────────────────────────────────────────────────────
\subsection{Key Innovations Beyond the Core Methods}
\label{subsec:key-innovations}

While the seven mathematical methods (M1--M7) constitute the
theoretical core of this dissertation, their integration into a
production SOC platform necessitated six additional innovation
categories, each representing a substantive engineering and research
contribution.

\paragraph{1. LLM Security Testing Suite.}
The growing deployment of large language models within security
pipelines introduces novel attack surfaces that traditional IDS
architectures do not address [148].  RobustIDPS.ai includes four
dedicated modules for LLM security evaluation.  The \emph{Prompt
Injection} module implements a live \texttt{DefensePipeline} that
processes user inputs through three sequential stages: (i)~regex-based
pattern matching against 8~injection categories (direct override,
context manipulation, role hijacking, encoding attacks, nested
injection, multi-turn exploitation, system prompt extraction, and
output formatting attacks), (ii)~boundary enforcement that validates
input--output scope constraints, and (iii)~output filtering that
sanitises model responses before delivery to the analyst.  The
\emph{Jailbreak Taxonomy} module catalogues 8~jailbreak technique
families---including DAN-style prompts, character roleplay, hypothetical
framing, token smuggling, multi-language bypass, instruction hierarchy
exploitation, context window overflow, and reward hacking---and
provides live testing interfaces against real LLM endpoints to
measure defence efficacy [149].  The \emph{RAG Poisoning} module
demonstrates attacks against Retrieval-Augmented Generation pipelines
through document provenance tracking, Jaccard similarity-based
retrieval scoring, and poisoning detection heuristics that flag
documents with anomalous similarity distributions.  Finally, the
\emph{Multi-Agent Chain} module models a 4-agent topology (Planner,
Executor, Validator, Reporter) with inter-agent trust scoring,
enabling analysts to observe how adversarial manipulation of one agent
can cascade through the chain [150].

\paragraph{2. SOC Intelligence Suite.}
Drawing inspiration from commercial platforms such as Microsoft
Security Copilot, SentinelOne Purple AI, and Google Sec-Gemini
v1~[151], the SOC Intelligence suite provides six analyst-facing
tools.  \emph{Auto-Investigation} implements a 4-phase agentic
pipeline: (a)~alert enrichment via contextual feature extraction,
(b)~hypothesis generation using LLM-assisted reasoning,
(c)~evidence correlation across temporal and topological dimensions,
and (d)~verdict synthesis with confidence calibration.
\emph{Natural-Language Threat Hunt} exposes a query parser that
translates analyst questions into structured filters over attack type,
confidence threshold, port range, and severity level, thereby
lowering the barrier to proactive hunting for analysts without
specialised query-language training.  \emph{Incident Reports} are
auto-generated with MITRE ATT\&CK coverage annotations, executive
summaries, and technical appendices.  \emph{Threat Intel} provides IP
reputation scoring with geolocation enrichment.  The \emph{Rule
Generator} offers 13~Suricata/Snort rule templates covering common
attack signatures, enabling one-click rule synthesis from detected
anomalies.  The \emph{CVE Mapper} links 6~attack-class-to-CVE
mappings with associated CVSS v3.1 base scores, providing analysts
with vulnerability context for observed threats [152].

\paragraph{3. Multi-Agent PQC-IDS.}
A novel contribution extending the PQ-IDPS work (Branch~3 of the
dissertation) is the \emph{Multi-Agent PQC-IDS}, a cooperative
4-agent model comprising 70{,}598 trainable parameters [153].  The
\emph{Traffic Analyst} agent maps 83~input features to 34~attack
classes using a two-layer feed-forward network with batch
normalisation.  The \emph{PQC Specialist} agent focuses on
post-quantum cryptographic traffic, classifying 83~features into
14~PQ algorithm categories.  The \emph{Anomaly Detector} agent
employs an autoencoder with architecture $83 \to 16 \to 83$,
flagging flows whose reconstruction error exceeds a learned
threshold.  The \emph{Coordinator} agent performs attention-based
fusion over the three specialist outputs, producing a unified
detection verdict.  The model is publicly available at
\texttt{github.com/rogerpanel/Multi-Agent-PQC-models}, and the
training dataset is archived at DOI:
10.34740/kaggle/dsv/15424420~[154], comprising subsets from
CESNET-TLS22 and ArielCyber PQClass.

\paragraph{4. Continual Learning and Autonomous Response.}
To address concept drift in production traffic streams, the platform
implements Elastic Weight Consolidation (EWC) with experience replay
(buffer size $M = 5{,}000$ flows) [155].  The autonomous response
module uses Constrained Policy Optimisation (CPO) with 5~graduated
response actions ordered by severity: \textsc{Monitor} $\to$
\textsc{Log-Enrich} $\to$ \textsc{Rate-Limit} $\to$
\textsc{Isolate} $\to$ \textsc{Quarantine}.  A unified Fisher
Information matrix with blending coefficient $\beta = 0.7$ governs
the trade-off between plasticity and stability, while a bounded
autonomy configuration ensures that high-severity actions
(\textsc{Isolate}, \textsc{Quarantine}) require human-in-the-loop
confirmation before execution [156].

\paragraph{5. MLSecOps Standards Compliance.}
The platform provides structured compliance tooling aligned with three
major frameworks.  The \emph{MITRE ATT\&CK Mapper} associates each of
the 30~attack classes detected by the model ensemble with
corresponding ATT\&CK techniques and sub-techniques, enabling
threat-intelligence interoperability.  An \emph{OWASP LLM Top~10
Scorecard} evaluates the platform's own LLM integrations against each
risk category (LLM01--LLM10), with pass/fail indicators and
remediation guidance.  A \emph{NIST AI RMF Dashboard} tracks
alignment with the AI Risk Management Framework across its four
functions (Govern, Map, Measure, Manage).  The \emph{Alert Triage}
classifier applies a severity-weighted prioritisation algorithm that
combines model confidence, attack class severity, and contextual
indicators to rank alerts for analyst review [157].

\paragraph{6. Security Hardening.}
Production deployment demands defence-in-depth beyond the detection
algorithms themselves.  All user-facing content is sanitised through
\textbf{DOMPurify} to prevent cross-site scripting (XSS) attacks
against the analyst interface.  WebSocket connections require
\textbf{JWT token validation} at the handshake phase, preventing
unauthorised subscription to live traffic streams.  A \textbf{per-job
authorisation} model ensures that background tasks (model inference,
report generation) execute only under the identity of the requesting
user.  \textbf{Security headers middleware} enforces
Content-Security-Policy, X-Frame-Options, and
Strict-Transport-Security across all HTTP responses [158].

% ────────────────────────────────────────────────────────────
\subsection{Validation Datasets and Methodology}
\label{subsec:validation-datasets}

Validation of the integrated platform requires purpose-built datasets
that exercise the full breadth of detection capabilities.  Three
crafted PCAP captures were generated for this purpose, each designed
to test a distinct operational scenario:

\begin{enumerate}
    \item \textbf{validation\_benchmark.pcap} (323.5\,MB): Contains
    60{,}000 synthetic flows spanning 34~attack classes, with a
    60/40 split between labelled and unlabelled flows.  The labelled
    subset provides ground-truth annotations for supervised accuracy
    measurement, while the unlabelled subset tests the platform's
    unsupervised anomaly detection capabilities (M2, M6) and
    semi-supervised graph methods (M1, M4).

    \item \textbf{pqc\_validation\_benchmark.pcap}: A PQC-focused
    capture containing 40{,}000 flows generated from 10~post-quantum
    cryptographic algorithm implementations (Kyber-512, Kyber-768,
    Kyber-1024, Dilithium-2, Dilithium-3, Dilithium-5, SPHINCS+-128f,
    Falcon-512, Falcon-1024, Classic McEliece) and 4~PQ-specific
    attack variants (downgrade, oracle, hybrid-mode exploitation, and
    side-channel simulation).

    \item \textbf{combined\_validation\_benchmark.pcap}: An integrated
    capture of 50{,}000 flows with a deliberately heterogeneous
    composition: 40\% conventional IDS traffic, 30\% PQC-specific
    traffic, and 30\% adversarially perturbed flows designed to test
    robustness under evasion conditions.
\end{enumerate}

These crafted datasets are complemented by external PQC datasets
sourced from Kaggle (DOI: 10.34740/kaggle/dsv/15424420) [154],
including subsets from CESNET-TLS22 (Czech national research network
TLS captures), ArielCyber PQClass (post-quantum traffic
classification), PQS TLS Measurements (PQ algorithm handshake
timings), and PQ IoT Impact (resource-constrained device
benchmarks).  Together, these sources provide coverage across both
conventional and post-quantum threat landscapes.

The validation methodology follows a four-stage pipeline: (i)~upload
of PCAP files via the Live Monitor File Replay interface, which
deterministically re-ingests packets at configurable rates;
(ii)~multi-model detection, where each flow is classified by all
applicable models in the registry and results are aggregated through
the tropical semiring fusion layer (M7); (iii)~cross-page analysis,
in which detection outputs are correlated across the Analytics,
Alert Causality, and Auto-Investigation pages to assess end-to-end
SOC workflow coverage; and (iv)~comparison against ground-truth CSV
annotations, computing standard metrics (accuracy, precision, recall,
F1-score, and area under the ROC curve) per attack class and per
model [159].  This methodology ensures that validation captures not
only individual model performance but also the integrative behaviour
of the platform as a cohesive analytical instrument.


% Continues from ch3_validation_seg1.tex
% Subsections 3.18.4 through 3.18.6 + Bibliography

\subsection{Platform Validation Results}
\label{subsec:platform-validation-results}

We present comprehensive validation results across three evaluation dimensions:
model classification performance, adversarial robustness, and LLM defense
pipeline effectiveness. All experiments were conducted on the CIC-IoT-2023
benchmark dataset comprising 34 attack classes, using a workstation equipped
with an NVIDIA A100 GPU (80\,GB VRAM) and 128\,GB system memory.

\subsubsection{Model Classification Performance}

Table~\ref{tab:model-performance} summarizes the detection performance of
all eight models integrated into the RobustIDPS.ai platform. Each model was
trained on an 80/10/10 train/validation/test split with five-fold
cross-validation, and we report mean metrics across folds.

\begin{table}[!htbp]
\centering
\caption{Model Performance on Validation Benchmark (CIC-IoT-2023, 34 Classes)}
\label{tab:model-performance}
\renewcommand{\arraystretch}{1.15}
\begin{tabular}{@{}lcccc@{}}
\toprule
\textbf{Model} & \textbf{Acc.} & \textbf{F1} & \textbf{FPR} & \textbf{Lat.\,(ms)} \\
\midrule
SurrogateIDS (7-Branch)      & 97.8\% & 0.976 & 1.4\% & 0.8 \\
CL-RL Unified + EWC          & 96.8\% & 0.965 & 1.8\% & 1.1 \\
CT-TGNN (Neural ODE)         & 96.2\% & 0.958 & 2.1\% & 1.2 \\
SDE-TGNN (Stochastic)        & 96.0\% & 0.955 & 2.3\% & 1.4 \\
MambaShield (SSM)             & 95.9\% & 0.954 & 2.2\% & 0.5 \\
Stochastic Transformer        & 95.4\% & 0.949 & 2.5\% & 1.8 \\
FedGTD (Byzantine-resilient)  & 95.1\% & 0.946 & 2.7\% & 1.3 \\
Multi-Agent PQC-IDS           & 94.3\% & 0.938 & 2.9\% & 1.5 \\
\bottomrule
\end{tabular}
\end{table}

The SurrogateIDS ensemble achieves the highest accuracy at 97.8\%, which is
expected given its seven-branch architecture that combines heterogeneous
classifiers through learned gating weights. Notably, MambaShield delivers
competitive accuracy (95.9\%) at the lowest inference latency (0.5\,ms),
making it particularly suitable for high-throughput edge deployments where
per-packet processing budgets are constrained. The Multi-Agent PQC-IDS model
exhibits lower baseline accuracy (94.3\%), which reflects the additional
complexity of classifying post-quantum cryptographic handshake patterns---a
task for which training data remains comparatively scarce [143].

\subsubsection{Adversarial Robustness Evaluation}

Table~\ref{tab:adversarial-robustness} presents the robustness evaluation
of the SurrogateIDS ensemble under six representative attack methods. We
report clean accuracy, attacked accuracy without defenses, and robust
accuracy after adversarial training (AT) with the projected gradient descent
(PGD) formulation [144].

\begin{table}[!htbp]
\centering
\caption{Adversarial Robustness Under Six Attack Methods}
\label{tab:adversarial-robustness}
\renewcommand{\arraystretch}{1.15}
\begin{tabular}{@{}lcccc@{}}
\toprule
\textbf{Attack} & \textbf{$\varepsilon/\sigma$} & \textbf{Clean} & \textbf{Attacked} & \textbf{Robust (AT)} \\
\midrule
FGSM             & $\varepsilon{=}0.1$ & 97.8\% & 89.2\% & 93.1\% \\
PGD-40           & $\varepsilon{=}0.1$ & 97.8\% & 84.7\% & 89.6\% \\
C\&W             & $\kappa{=}10$       & 97.8\% & 82.3\% & 85.2\% \\
DeepFool         & 30 iter             & 97.8\% & 86.1\% & 90.8\% \\
Gaussian Noise   & $\sigma{=}0.1$      & 97.8\% & 91.4\% & 95.3\% \\
Label Masking    & 10\% flip           & 97.8\% & 88.9\% & 93.7\% \\
\bottomrule
\end{tabular}
\end{table}

The C\&W attack proves most effective at degrading classifier accuracy
(82.3\% under attack), consistent with its optimization-based formulation
that directly minimizes perturbation magnitude. Adversarial training
recovers between 2.9 and 5.3 percentage points across all attack types,
with the strongest recovery observed against FGSM (+3.9\,pp) and Gaussian
noise (+3.9\,pp). The relatively modest recovery against C\&W (+2.9\,pp)
underscores the difficulty of defending against optimization-based attacks
and motivates continued research into certified robustness guarantees [145].

\subsubsection{LLM Defense Pipeline Evaluation}

Table~\ref{tab:llm-defense} reports detection effectiveness of the
multi-layer LLM defense pipeline across eight prompt injection categories
derived from the OWASP Top~10 for LLM Applications taxonomy [146].

\begin{table}[!htbp]
\centering
\caption{LLM Defense Pipeline Detection Effectiveness}
\label{tab:llm-defense}
\renewcommand{\arraystretch}{1.15}
\begin{tabular}{@{}lccc@{}}
\toprule
\textbf{Injection Category} & \textbf{Det.\,Rate} & \textbf{Bypass} & \textbf{Conf.} \\
\midrule
Direct Override          & 98.5\% & 1.5\%  & 0.96 \\
Context Manipulation     & 94.2\% & 5.8\%  & 0.89 \\
Role-Playing             & 91.7\% & 8.3\%  & 0.84 \\
Encoding/Obfuscation     & 96.3\% & 3.7\%  & 0.91 \\
Multi-Turn Conversation  & 88.4\% & 11.6\% & 0.78 \\
System Prompt Extraction & 97.1\% & 2.9\%  & 0.93 \\
Tool/Function Manip.     & 93.8\% & 6.2\%  & 0.87 \\
Data Exfiltration        & 95.6\% & 4.4\%  & 0.90 \\
\midrule
\textbf{Overall (8 cat.)} & \textbf{94.5\%} & \textbf{5.5\%} & \textbf{0.89} \\
\bottomrule
\end{tabular}
\end{table}

The defense pipeline achieves an aggregate detection rate of 94.5\% across
all injection categories, with direct override attempts being the most
reliably detected (98.5\%) due to their distinctive lexical signatures.
Multi-turn conversation attacks present the greatest challenge (88.4\%
detection), as they distribute malicious intent across multiple exchanges,
making single-turn pattern matching insufficient. This finding aligns with
observations from recent commercial LLM security benchmarks [147] and
highlights the importance of the stateful context tracking implemented in
our pipeline. For operational deployment, we recommend coupling the
automated pipeline with human-in-the-loop review for low-confidence
detections (confidence $<0.80$), which would affect approximately 12\% of
flagged interactions based on our validation data.

\subsection{Contributions to the Research Community}
\label{subsec:contributions}

The RobustIDPS.ai platform and its associated artifacts are released as
open-source contributions to facilitate reproducibility and accelerate
further research in adversarially robust intrusion detection.

\textbf{Software Release.} The complete platform source code, including all
eight model implementations, the LLM defense pipeline, and the interactive
dashboard, is archived at Zenodo with DOI~\texttt{10.5281/zenodo.19129512}
[148]. The repository is maintained under the MIT license at
\url{https://github.com/rogerpanel/robustidps.ai}.

\textbf{Multi-Agent PQC-IDS Models.} Pre-trained model weights and training
scripts for the Multi-Agent PQC-IDS architecture are released separately at
\url{https://github.com/rogerpanel/Multi-Agent-PQC-models}, enabling
researchers to reproduce our post-quantum classification results without
retraining from scratch.

\textbf{PQC Benchmark Datasets.} Curated datasets for post-quantum
cryptographic traffic classification are published on Kaggle with DOI
\texttt{10.34740/kaggle/dsv/15424420} [149], containing labeled PCAP
captures for CRYSTALS-Kyber, CRYSTALS-Dilithium, SPHINCS+, and hybrid
TLS~1.3 handshakes.

\textbf{Validation PCAP Generators.} Three reproducible benchmark generation
scripts---\texttt{generate\_validation\_pcap.py},
\texttt{generate\_pqc\_validation\_pcap.py}, and
\texttt{generate\_combined\_validation\_pcap.py}---are included in the
\texttt{sample\_data/} directory, enabling deterministic recreation of our
evaluation datasets with configurable attack distributions and traffic
volumes.

\textbf{Publications.} Two IEEE TNNLS draft papers are currently in
preparation: one addressing adversarially robust continual learning for IDS,
and a second on securing LLM-integrated network defense systems.

\textbf{Citation.} Researchers using this platform should cite:
Anaedevha,~R.\,N. and Trofimov,~A.\,G. (2026). \textit{RobustIDPS.ai}
(Version~1.1.0) [Computer software]. Zenodo.
\url{https://doi.org/10.5281/zenodo.19129512}.

\subsection{Future Research Directions}
\label{subsec:future-directions}

While the RobustIDPS.ai platform demonstrates strong validation results
across classification, robustness, and LLM security dimensions, several
promising research directions remain open.

\textbf{1.\;Real Multi-Agent LLM Orchestration.} The current implementation
employs pattern-based dispatch to route network alerts to specialized
analysis modules. A natural extension is to replace this with genuine
multi-agent LLM orchestration, where each agent (triage, forensics,
response planning, compliance) is backed by a dedicated language model
instance capable of autonomous reasoning and inter-agent communication.
This would enable richer contextual analysis but introduces challenges
around coordination latency, hallucination propagation, and cost management
that must be addressed before production deployment [150].

\textbf{2.\;Online Continual Learning with Active Forgetting.} Our current
continual learning pipeline operates in batch mode with periodic
retraining. Extending this to true online learning with streaming data
ingestion would reduce adaptation latency from hours to minutes. Equally
important is the incorporation of active forgetting mechanisms that
selectively discard obsolete attack signatures, preventing unbounded memory
growth while preserving detection capability for recurrent threat patterns.

\textbf{3.\;Multi-Modal Sensor Fusion.} The platform currently processes
network traffic features in isolation. Fusing network telemetry with
complementary data sources---system call logs, endpoint detection and
response (EDR) telemetry, and application-layer audit trails---would provide
a more holistic view of attack campaigns, particularly for lateral movement
and privilege escalation scenarios that manifest across multiple
observability layers [151].

\textbf{4.\;Post-Quantum IDS on Constrained IoT Devices.} Deploying
PQC-aware intrusion detection on resource-constrained IoT devices remains
an open challenge. Leveraging the PQ IoT Impact dataset [152] and
techniques such as model quantization, knowledge distillation, and
neural architecture search could yield compact classifiers suitable for
microcontroller-class hardware with sub-megabyte memory budgets.

\textbf{5.\;Federated PQC Model Training with Differential Privacy.}
Extending the FedGTD framework to support cross-organizational PQC model
training with formal differential privacy guarantees ($\varepsilon$-DP)
would enable collaborative threat intelligence sharing without exposing
sensitive network topologies or traffic patterns between participating
organizations.

\textbf{6.\;Formal Verification of LLM Defense Pipelines.} Current prompt
injection defenses rely on empirical evaluation against known attack
taxonomies. Applying formal verification techniques---such as abstract
interpretation or satisfiability modulo theories (SMT) solvers---to prove
injection detection completeness over defined input grammars would provide
stronger security guarantees than empirical testing alone.

\textbf{7.\;Cross-Platform Benchmarking Against Commercial SOC Platforms.}
A rigorous comparative evaluation against production-grade AI security
platforms---including CrowdStrike Charlotte~AI [143], Microsoft Security
Copilot [144], SentinelOne Purple~AI [145], and Google
Sec-Gemini [146]---would contextualize our academic results within the
operational landscape and identify specific capability gaps requiring
further development.

% === BIBLIOGRAPHY ===

\begin{thebibliography}{99}

\bibitem{crowdstrike_charlotte}
CrowdStrike, ``Charlotte AI: Autonomous AI Security Analyst,'' 2025.
[Online]. Available:
\url{https://www.crowdstrike.com/en-us/platform/charlotte-ai/}

\bibitem{ms_copilot}
Microsoft, ``Microsoft Security Copilot,'' 2025. [Online]. Available:
\url{https://learn.microsoft.com/en-us/copilot/security/}

\bibitem{sentinelone_purple}
SentinelOne, ``Purple AI: Autonomous SecOps Analyst,'' 2025. [Online].
Available: \url{https://www.sentinelone.com/platform/purple/}

\bibitem{sec_gemini}
Google, ``Sec-Gemini v1: Experimental Cybersecurity AI Model,'' 2025.
[Online]. Available:
\url{https://security.googleblog.com/2025/04/google-launches-sec-gemini-v1-new.html}

\bibitem{omnisec}
C.~Chen \textit{et al.}, ``OMNISEC: LLM-Driven Provenance-based Intrusion
Detection via Retrieval-Augmented Behavior Prompting,''
\textit{arXiv:2503.03108}, 2025.

\bibitem{owasp_llm}
OWASP, ``OWASP Top 10 for Large Language Model Applications,'' Version~1.1,
2024.

\bibitem{nist_ai_rmf}
NIST, ``Artificial Intelligence Risk Management Framework (AI RMF 1.0),''
NIST AI 100-1, 2023.

\bibitem{pqclass}
ArielCyber, ``PQClass: Classifying Post-Quantum Cryptography in Encrypted
Network Traffic,'' \textit{IEEE ICC}, 2025.

\bibitem{pqs_tls}
M.~Sosnowski \textit{et al.}, ``The Performance of Post-Quantum TLS~1.3,''
in \textit{Proc. CoNEXT}, 2023.

\bibitem{pqc_iot}
L.~Cruz-Piris \textit{et al.}, ``Measuring the impact of post-quantum
cryptography in Industrial IoT scenarios,'' \textit{Internet of Things},
vol.~34, 2025.

\bibitem{hybrid_llm_ids}
``Hybrid LLM-Enhanced Intrusion Detection for Zero-Day Threats in IoT
Networks,'' \textit{arXiv:2507.07413}, 2025.

\end{thebibliography}

\end{document}
